\chapter{函数与变化}\label{ch:g11:func}

本章梳理学习微积分之前关于\omterm{def:g11:func:function}{函数}需要掌握的知识：
一小类\omterm{def:g10:functions:graph}{图像}须熟记的\emph{\omterm{def:g11:func:reference}{基本函数}}，
并精确给出\emph{\omterm{def:g11:func:monotone}{递增}}与
\emph{\omterm{def:g10:functions:variations}{递减}}的确切含义，以及\omterm{def:g11:func:function}{函数}组合后变化的规律。

\section{函数与图像}

\begin{definition}[函数、定义域、图像]\label{def:g11:func:function}
称 $f$ 为\emph{函数}\index{函数}，它对集合中的每个数 $x$
$D \subseteq \R$（其\emph{定义域}\index{定义域}）都恰好指定一个数
$f(x)$。它的\emph{\omterm{def:g10:functions:graph}{图像}}是点集
$\bigl(x, f(x)\bigr)$，其中 $x \in D$。
\end{definition}

\begin{example}
\omterm{def:g11:func:function}{函数} $f(x) = \frac{\sqrt{x+2}}{x-1}$ 要求 $x + 2 \geq 0$（根号内非负）
且 $x \neq 1$（分母不为零），故\omterm{def:g11:func:function}{定义域}为
$\intco{-2}{1} \cup \intoo{1}{+\infty}$。
\end{example}

\section{基本函数}

下列五个\omterm{def:g11:func:function}{函数}及其\omterm{def:g10:functions:graph}{图像}必须烂熟于心。

\begin{definition}[基本函数]\label{def:g11:func:reference}
\index{基本函数}
\[
x \mapsto x^2 \text{（平方）}, \quad
x \mapsto x^3 \text{（立方）}, \quad
x \mapsto \sqrt{x}, \quad
x \mapsto \frac{1}{x}, \quad
x \mapsto \abs{x}.
\]
它们的\omterm{def:g11:func:function}{定义域}依次为 $\R$、$\R$、$\intco{0}{+\infty}$、$\R \setminus \{0\}$
和 $\R$。$x \mapsto \frac1x$ 的\omterm{def:g10:functions:graph}{图像}是
一条\emph{双曲线}\index{双曲线}。
\end{definition}

\begin{omfigure}
\begin{tikzpicture}
\begin{axis}[omaxis,
  width=7cm, height=5.6cm,
  xmin=-2.3, xmax=2.3, ymin=-3.4, ymax=3.6,
  xtick={-1,1}, ytick={-1,1}]
  \addplot[omDef, thick, domain=-1.85:1.85] {x^2};
  \addplot[omThm, thick, domain=-1.5:1.5] {x^3};
  \node[omDef, font=\small, anchor=east] at (axis cs:-1.5,1.9) {$x^2$};
  \node[omThm, font=\small, anchor=east] at (axis cs:-1.32,-1.73) {$x^3$};
\end{axis}
\end{tikzpicture}%
\begin{tikzpicture}
\begin{axis}[omaxis,
  width=7cm, height=5.6cm,
  xmin=-3.4, xmax=4.4, ymin=-3.4, ymax=3.6,
  xtick={1}, ytick={1}]
  \addplot[omProp, thick, domain=-3.2:-0.31, samples=60] {1/x};
  \addplot[omProp, thick, domain=0.31:4.2, samples=60] {1/x};
  \addplot[omDef, thick, domain=0:4.2] {sqrt(x)};
  \addplot[omThm, thick, domain=-3.2:3.2] {abs(x)};
  \node[omDef, font=\small, anchor=north] at (axis cs:3.6,1.55)
    {$\sqrt{x}$};
  \node[omThm, font=\small, anchor=east] at (axis cs:-2.4,2.8)
    {$\abs{x}$};
  \node[omProp, font=\small, anchor=west] at (axis cs:0.75,2.4)
    {$\tfrac1x$};
\end{axis}
\end{tikzpicture}

{\small 五条\omterm{def:g11:func:reference}{基本函数}的\omterm{def:g10:functions:graph}{图像}：\omterm{def:g11:quad:parabola}{抛物线} $x^2$ 与立方\omterm{def:g11:func:function}{函数} $x^3$
（左图）；$\sqrt x$、$\abs x$ 与\omterm{def:g11:func:reference}{双曲线} $\frac1x$（右图）。}
\end{omfigure}

\section{区间上的变化}

\begin{definition}[递增、递减]\label{def:g11:func:monotone}
设 $f$ 在\omterm{def:g10:numbers:interval}{区间} $I$ 上有定义。称 $f$ 在 $I$ 上
\emph{递增}\index{递增函数}，若它保持次序：
\[
\text{对所有 } u, v \in I: \quad u < v \implies f(u) \leq f(v),
\]
而若它颠倒次序，则称其\emph{\omterm{def:g10:functions:variations}{递减}}（$u < v \implies f(u) \geq
f(v)$）。若右边的不等式为严格的，则称 $f$ 为\emph{严格}
递增（或\omterm{def:g10:functions:variations}{递减}）。凡在 $I$ 上递增或\omterm{def:g10:functions:variations}{递减}的\omterm{def:g11:func:function}{函数}，
称为 $I$ 上的\emph{单调}\omterm{def:g11:func:function}{函数}。
\end{definition}

\begin{method}[依定义证明变化]\label{met:g11:func:variation}
在\omterm{def:g10:numbers:interval}{区间}内任取两数 $u < v$，确定差 $f(v) - f(u)$
的符号，通常将其\emph{\omterm{def:g10:algebra:expand}{因式分解}}。
若该符号在\omterm{def:g10:numbers:interval}{区间}上恒定，即可下结论。
\end{method}

\begin{proposition}[基本函数的变化]\label{prop:g11:func:refvariations}
\begin{enumerate}
  \item $x^2$ 在 $\intoc{-\infty}{0}$ 上严格\omterm{def:g10:functions:variations}{递减}，且
        在 $\intco{0}{+\infty}$ 上严格\omterm{def:g11:func:monotone}{递增}。
  \item $x^3$ 在 $\R$ 上严格\omterm{def:g11:func:monotone}{递增}。
  \item $\sqrt x$ 在 $\intco{0}{+\infty}$ 上严格\omterm{def:g11:func:monotone}{递增}。
  \item $\frac1x$ 在 $\intoo{-\infty}{0}$ 上严格\omterm{def:g10:functions:variations}{递减}，且在
        $\intoo{0}{+\infty}$ 上亦严格\omterm{def:g10:functions:variations}{递减}（但\emph{不}在它们的\omterm{def:g10:numbers:interunion}{并集}上）。
  \item $\abs x$ 在 $\intoc{-\infty}{0}$ 上严格\omterm{def:g10:functions:variations}{递减}，且
        在 $\intco{0}{+\infty}$ 上严格\omterm{def:g11:func:monotone}{递增}。
\end{enumerate}
\end{proposition}

\begin{proof}
设 $u < v$。在每种情形下对 $f(v) - f(u)$ 作\omterm{def:g10:algebra:expand}{因式分解}。

\emph{1.} $v^2 - u^2 = (v-u)(v+u)$。在 $\intco{0}{+\infty}$ 上：$v - u > 0$
且 $v + u > 0$（因 $v > u \geq 0$），故 $v^2 > u^2$。在
$\intoc{-\infty}{0}$ 上：$v + u < 0$，故 $v^2 < u^2$。

\emph{2.} $v^3 - u^3 = (v - u)(v^2 + uv + u^2)$，其中第二个因式
等于 $\left(v + \frac u2\right)^2 + \frac{3u^2}{4} > 0$，除非
$u = v = 0$。由此可得 $v^3 > u^3$（当 $v > u$ 时）。

\emph{3.} 当 $0 \leq u < v$ 时：
$\sqrt v - \sqrt u = \frac{(\sqrt v - \sqrt u)(\sqrt v + \sqrt u)}
{\sqrt v + \sqrt u} = \frac{v - u}{\sqrt v + \sqrt u} > 0$.

\emph{4.} 当 $0 < u < v$ 时：
$\frac1v - \frac1u = \frac{u - v}{uv} < 0$，因 $u - v < 0$ 且 $uv > 0$；
对 $u < v < 0$ 作同样的计算。在\omterm{def:g10:numbers:interunion}{并集}上结论不成立：
$-1 < 1$，但 $\frac{1}{-1} = -1 < 1 = \frac11$。

\emph{5.} 在 $\intco{0}{+\infty}$ 上，$\abs x = x$ 严格\omterm{def:g11:func:monotone}{递增}；在
$\intoc{-\infty}{0}$ 上，$\abs x = -x$ 严格\omterm{def:g10:functions:variations}{递减}。
\end{proof}

\begin{remark}
第 4 点是一个经典陷阱：“在 $\intoo{-\infty}{0}$ 上\omterm{def:g10:functions:variations}{递减}，且在
$\intoo{0}{+\infty}$ 上\omterm{def:g10:functions:variations}{递减}”\emph{不}等同于“在
$\R \setminus \{0\}$ 上\omterm{def:g10:functions:variations}{递减}”。单调性只在\omterm{def:g10:numbers:interval}{区间}上才有意义。
\end{remark}

\begin{example}[不计算而比较大小]
由于 $\sqrt x$ 严格\omterm{def:g11:func:monotone}{递增}，$\sqrt{7} < \sqrt{8}$；由于
$x^2$ 在 $\intoc{-\infty}{0}$ 上严格\omterm{def:g10:functions:variations}{递减}，
故 $(-3)^2 > (-2)^2$。单调性把像的比较
转化为自变量的比较。
\end{example}

\section{由旧函数构造新函数}

\begin{proposition}[变换]\label{prop:g11:func:transforms}
设 $u$ 为定义在\omterm{def:g10:numbers:interval}{区间} $I$ 上的\omterm{def:g11:func:function}{函数}，$k \in \R$。
\begin{enumerate}
  \item $u + k$ 与 $u$ 在 $I$ 上的单调性\emph{相同}；其\omterm{def:g10:functions:graph}{图像}是将
        $u$ 的\omterm{def:g10:functions:graph}{图像}沿竖直方向平移 $k$ 个单位所得。
  \item 当 $\lambda > 0$ 时，$\lambda u$ 与 $u$ 的单调性相同；当
        $\lambda < 0$ 时，单调性相反。
  \item 若 $u \geq 0$ 在 $I$ 上成立，则 $\sqrt u$ 与
        $u$ 的单调性相同。
  \item 若 $u \neq 0$ 且 $u$ 在 $I$ 上保持定号，则 $\frac1u$ 的
        单调性与 $u$ 的\emph{相反}。
\end{enumerate}
\end{proposition}

\begin{proof}
设 $u < v$ 为 $I$ 中的两点（记号略有滥用，我们比较的是 $x$ 值 $u < v$；记
$s = u(x_1)$，$t = u(x_2)$ 分别是 $x_1 < x_2$ 的像）。

\emph{1.} $(u+k)(x_2) - (u+k)(x_1) = t - s$：像之差
不变，因此其符号也不变。

\emph{2.} $\lambda t - \lambda s = \lambda(t - s)$：当
$\lambda > 0$ 时符号保持，$\lambda < 0$ 时符号改变。

\emph{3.} 若 $u$ \omterm{def:g11:func:monotone}{递增}且 $x_1 < x_2$，则 $0 \leq s \leq t$，
由 $\sqrt{}$ \omterm{def:g11:func:monotone}{递增}（见\cref{prop:g11:func:refvariations}）可得
$\sqrt s \leq \sqrt t$。\omterm{def:g10:functions:variations}{递减}的情形同理。

\emph{4.} 若 $u$ \omterm{def:g11:func:monotone}{递增}，$0 < s \leq t$（或 $s \leq t < 0$），则
$\frac1s \geq \frac1t$，因为 $x \mapsto \frac1x$ 在
$0$ 两侧各自\omterm{def:g10:functions:variations}{递减}，次序反转。
\end{proof}

\begin{example}\label{ex:g11:func:composite}
研究\omterm{def:g11:func:function}{函数} $f(x) = \sqrt{x^2 + 1}$ 在 $\R$ 上的单调性。内层\omterm{def:g11:func:function}{函数}
$u(x) = x^2 + 1$ 为正，在 $\intoc{-\infty}{0}$ 上\omterm{def:g10:functions:variations}{递减}，且
在 $\intco{0}{+\infty}$ 上\omterm{def:g11:func:monotone}{递增}（即平移后的平方\omterm{def:g11:func:function}{函数}）。由第 3 点，$f$
的单调性相同：先\omterm{def:g10:functions:variations}{递减}后\omterm{def:g11:func:monotone}{递增}，\omterm{def:g10:functions:extrema}{最小值}为
$f(0) = 1$。
\end{example}

\begin{omfigure}
\begin{tikzpicture}
\begin{axis}[omaxis,
  width=8.6cm, height=5.8cm,
  xmin=-2.6, xmax=2.6, ymin=-1.7, ymax=4.6,
  xtick={-1,1}, ytick={1,2}]
  \addplot[omDef, thick, domain=-1.9:1.9] {x^2};
  \addplot[omThm, thick, domain=-1.55:1.55] {x^2 + 1};
  \addplot[omProp, thick, domain=-2.4:2.4, samples=80]
    {sqrt(x^2 + 1)};
  \node[omDef, font=\small, anchor=west] at (axis cs:1.9,3.4) {$x^2$};
  \node[omThm, font=\small, anchor=west] at (axis cs:1.35,3.3)
    {$x^2{+}1$};
  \node[omProp, font=\small, anchor=north] at (axis cs:-1.9,1.95)
    {$\sqrt{x^2+1}$};
\end{axis}
\end{tikzpicture}

{\small 变换示意：将 $x^2$ 向上平移 $1$ 个单位（红色）保持
其单调性；取平方根（橙色）后单调性再次保持，如
\cref{ex:g11:func:composite}所示。}
\end{omfigure}

\section{偶函数与奇函数}

\begin{definition}[奇偶性]\label{def:g11:func:parity}
设 $f$ 为\omterm{def:g11:func:function}{定义域} $D$ 上的\omterm{def:g11:func:function}{函数}，该\omterm{def:g11:func:function}{定义域}关于 $0$ 对称（即 $-x \in D$
当 $x \in D$ 时成立）。称 $f$ 为\emph{偶函数}\index{偶函数}，若
$f(-x) = f(x)$ 对所有 $x \in D$ 成立，另称为\emph{奇函数}\index{奇函数}，若
$f(-x) = -f(x)$ 对所有 $x \in D$ 成立。
\end{definition}

\begin{proposition}[几何意义]\label{prop:g11:func:paritysym}
\omterm{def:g11:func:parity}{偶函数}的\omterm{def:g10:functions:graph}{图像}关于 $y$ 轴对称；\omterm{def:g11:func:parity}{奇函数}的
\omterm{def:g10:functions:graph}{图像}关于原点对称。
\end{proposition}

\begin{proof}
若 $f$ 为\omterm{def:g11:func:parity}{偶函数}，且 $(x, y)$ 在曲线上（即 $y = f(x)$），则
$f(-x) = f(x) = y$，点 $(-x, y)$ 也在曲线上。若
$f$ 为\omterm{def:g11:func:parity}{奇函数}，$f(-x) = -y$，则 $(-x, -y)$ 为 $(x,y)$ 关于原点的对称点，
也在曲线上。
\end{proof}

\begin{example}
$x^2$ 与 $\abs x$ 是\omterm{def:g11:func:parity}{偶函数}；$x^3$ 与 $\frac1x$ 是\omterm{def:g11:func:parity}{奇函数}；
$f(x) = x^2 + x$ 既非\omterm{def:g11:func:parity}{奇函数}也非\omterm{def:g11:func:parity}{偶函数}：$f(-1) = 0$，而 $f(1) = 2$，故
$f(-1) \neq \pm f(1)$。利用奇偶性可减少一半的工作量：\omterm{def:g11:func:parity}{偶函数}或\omterm{def:g11:func:parity}{奇函数}
只需在 $\intco{0}{+\infty}$ 上研究。
\end{example}

\section{练习}

\begin{exercise}[$\star$]\label{exo:g11:func:1}
求下列各\omterm{def:g11:func:function}{函数}的\omterm{def:g11:func:function}{定义域}：
\[
f(x) = \sqrt{3 - x}, \qquad
g(x) = \frac{1}{x^2 - 4}, \qquad
h(x) = \frac{\sqrt{x}}{x - 5} .
\]
\end{exercise}

\begin{exercise}[$\star$]\label{exo:g11:func:2}
利用\omterm{def:g11:func:reference}{基本函数}的单调性，
不使用计算器比较下列各组数的大小：
\[
\sqrt{11} \text{ 与 } \sqrt{13}; \qquad
(-2.1)^2 \text{ 与 } (-2.05)^2; \qquad
\frac{1}{0.9} \text{ 与 } \frac{1}{1.1}; \qquad
(-1.01)^3 \text{ 与 } (-0.99)^3 .
\]
\end{exercise}

\begin{exercise}[$\star$]\label{exo:g11:func:3}
判断下列各\omterm{def:g11:func:function}{函数}是\omterm{def:g11:func:parity}{奇函数}、\omterm{def:g11:func:parity}{偶函数}，还是既非\omterm{def:g11:func:parity}{奇函数}也非\omterm{def:g11:func:parity}{偶函数}：
\[
f(x) = 3x^4 - x^2, \qquad
g(x) = x^3 + x, \qquad
h(x) = x^2 + x^3, \qquad
k(x) = \frac{x}{x^2+1} .
\]
\end{exercise}

\begin{exercise}[$\star$]\label{exo:g11:func:4}
\omterm{def:g11:func:function}{函数} $f$ 的单调性如下：在
$\intcc{-3}{1}$ 上从 $f(-3) = -2$ \omterm{def:g11:func:monotone}{递增}到 $f(1) = 4$，然后在
$\intcc{1}{5}$ 上\omterm{def:g10:functions:variations}{递减}到 $f(5) = 0$。\omterm{def:g10:algebra:equation}{方程}
$f(x) = 1$ 在 $\intcc{-3}{5}$ 上有多少个解？$f(x) = 5$ 呢？
\end{exercise}

\begin{exercise}[$\star\star$]\label{exo:g11:func:5}
利用\cref{met:g11:func:variation}（$f(v) - f(u)$ 的符号），证明
$f(x) = x^2 - 4x$ 在 $\intco{2}{+\infty}$ 上严格\omterm{def:g11:func:monotone}{递增}。
\end{exercise}

\begin{exercise}[$\star\star$]\label{exo:g11:func:6}
不计算任何导数，求下列\omterm{def:g11:func:function}{函数}的单调性：
\[
f(x) = \sqrt{5 - x} \ \text{ 于 } \intoc{-\infty}{5},
\quad
g(x) = \frac{1}{x^2 + 1} \ \text{ 于 } \intco{0}{+\infty},
\quad
h(x) = 3 - 2\sqrt{x} \ \text{ 于 } \intco{0}{+\infty}.
\]
\end{exercise}

\begin{exercise}[$\star\star$]\label{exo:g11:func:7}
设 $f(x) = (x-2)^2 - 3$，\omterm{def:g11:func:function}{定义域}为 $\R$。
\begin{enumerate}
  \item 求 $f$ 的单调性及其\omterm{def:g10:functions:extrema}{最小值}。
  \item 由此推出 $g(x) = \frac{1}{(x-2)^2+1}$ 的单调性及其
        \omterm{def:g10:functions:extrema}{最大值}。（注意 $g = \frac{1}{f + 4}$。）
\end{enumerate}
\end{exercise}

\begin{exercise}[$\star\star$]\label{exo:g11:func:8}
\omterm{def:g11:func:function}{函数} $f$ 的曲线经过点 $(0, 1)$、$(2, 3)$
及 $(4, 1)$，且 $f$ 在 $\intcc{0}{2}$ 上\omterm{def:g11:func:monotone}{递增}，在
$\intcc{2}{4}$ 上\omterm{def:g10:functions:variations}{递减}。画出一条可能的曲线，再画出
$f + 2$、$-f$ 及 $2f$ 的曲线。
\end{exercise}

\begin{exercise}[$\star\star$]\label{exo:g11:func:9}
证明\omterm{def:g11:func:function}{函数} $f(x) = \frac{2x+1}{x+1}$ 可以写成
$f(x) = 2 - \frac{1}{x+1}$，并由此推出其在
$\intoo{-1}{+\infty}$ 上的单调性。
\end{exercise}

\begin{exercise}[$\star\star$]\label{exo:g11:func:10}
判断下列命题的真假（说明理由或举出反例）：
\begin{enumerate}
  \item 若两个\omterm{def:g11:func:function}{函数}在 $I$ 上\omterm{def:g11:func:monotone}{递增}，则它们的和在 $I$ 上\omterm{def:g11:func:monotone}{递增}。
  \item 若两个\omterm{def:g11:func:function}{函数}在 $I$ 上\omterm{def:g11:func:monotone}{递增}，则它们的积在
        $I$ 上\omterm{def:g11:func:monotone}{递增}。
  \item 若 $f$ 在 $I$ 上\omterm{def:g11:func:monotone}{递增}，则 $f^2$（$= f \times f$）在
        $I$ 上\omterm{def:g11:func:monotone}{递增}。
\end{enumerate}
\end{exercise}

\begin{exercise}[$\star\star\star$]\label{exo:g11:func:11}
设 $f(x) = x + \frac{1}{x}$，$x > 0$。
\begin{enumerate}
  \item 证明：对于 $0 < u < v$，
        $f(v) - f(u) = (v - u)\,\frac{uv - 1}{uv}$.
  \item 由此推出 $f$ 在 $\intoc{0}{1}$ 上严格\omterm{def:g10:functions:variations}{递减}，
        在 $\intco{1}{+\infty}$ 上严格\omterm{def:g11:func:monotone}{递增}。
  \item 由此得出：$x + \frac1x \geq 2$ 对所有 $x > 0$ 成立，且
        仅当 $x = 1$ 时等号成立。
\end{enumerate}
\end{exercise}

\section{周末问题：对称性工厂}

\begin{problem}[{周末问题——每个函数都可以写成偶部分与奇部分之和：
奇偶性规则、变换游戏，以及
不用微积分求得的最小值}]
\label{pb:g11:func:1}
有些\omterm{def:g11:func:function}{函数}关于纵轴对称，有些关于原点中心对称，
多数两者皆非——然而一个小小的代数奇迹断言，
\emph{每个}\omterm{def:g11:func:function}{函数}都恰以一种方式，写成一个关于纵轴对称的部分
与一个关于原点中心对称的部分之和。本题将证明这一奇迹，
并在基本曲线上玩变换游戏
（见\cref{prop:g11:func:transforms}），最后以
差分\omterm{def:g10:algebra:expand}{因式分解}的技巧收尾，该技巧出现在
\cref{exo:g11:func:11}中，它能在导数出场前整整一章的
篇幅里解决最小化问题。

\medskip
\noindent\textbf{第一部分——奇偶性的熟练运用}
\begin{enumerate}
  \item 判断下列\omterm{def:g11:func:function}{函数}是\omterm{def:g11:func:parity}{偶函数}、\omterm{def:g11:func:parity}{奇函数}，还是两者皆非（见\cref{def:g11:func:parity}）？
        并逐一说明理由：$x^4 - 3x^2$；\ $x^3 - x$；\ $x^2 + x$；\
        $\abs x$；\ $\frac1x$。
  \item 说明每个结论的几何意义
        （见\cref{prop:g11:func:paritysym}）：\omterm{def:g11:func:parity}{偶函数}的\omterm{def:g10:functions:graph}{图像}具有何种对称性？
        \omterm{def:g11:func:parity}{奇函数}的\omterm{def:g10:functions:graph}{图像}呢？
  \item 奇偶性的乘法与正负号的乘法规则相同：证明两个\omterm{def:g11:func:parity}{奇函数}之积
        为\omterm{def:g11:func:parity}{偶函数}，并补全整个乘法表
        （偶 $\times$ 偶，偶 $\times$
        奇，奇 $\times$ 奇）。
  \item 直接推论：若 $f$ 为\omterm{def:g11:func:parity}{奇函数}且在 $0$ 处有定义，
        则 $f(0)$ 为多少？若 $f$ 为\omterm{def:g11:func:parity}{偶函数}且 $f(3) = 7$，则
        $f(-3)$ 为多少？并证明：只有零\omterm{def:g11:func:function}{函数}
        才既是\omterm{def:g11:func:parity}{偶函数}又是\omterm{def:g11:func:parity}{奇函数}。
  \item \omterm{def:g11:func:parity}{偶函数}的\omterm{def:g10:functions:graph}{图像}有多大一部分，就足以确定它的全部？
        给出一个既非\omterm{def:g11:func:parity}{偶函数}也非\omterm{def:g11:func:parity}{奇函数}的例子，并
        说明哪一条性质不成立。
\end{enumerate}

\medskip
\noindent\textbf{第二部分——分解定理}
设\omterm{def:g11:func:function}{函数} $f$ 的\omterm{def:g11:func:function}{定义域}关于 $0$ 对称，令
\[
E(x) = \frac{f(x) + f(-x)}{2},
\qquad
O(x) = \frac{f(x) - f(-x)}{2} .
\]
\begin{enumerate}[resume]
  \item 计算 $E$ 与 $O$，其中 $f(x) = x^3 + x^2 + 1$，并验证：
        $E$ 为\omterm{def:g11:func:parity}{偶函数}，$O$ 为\omterm{def:g11:func:parity}{奇函数}，且 $E + O = f$。
  \item 证明一般结论：对任意 $f$，\omterm{def:g11:func:function}{函数}
        $E$ 为\omterm{def:g11:func:parity}{偶函数}，$O$ 为\omterm{def:g11:func:parity}{奇函数}，且 $f = E + O$——\emph{每个
        \omterm{def:g11:func:function}{函数}都可分解为一个偶部分与一个奇部分}。
  \item 证明这种分解是唯一的：若
        $f = E_1 + O_1 = E_2 + O_2$，其中 $E_i$ 为\omterm{def:g11:func:parity}{偶函数}，$O_i$
        为\omterm{def:g11:func:parity}{奇函数}，对 $E_1 - E_2$ 运用第 4 题的结论。
  \item 分解 $f(x) = \dfrac{1}{1 - x}$（其中
        $x \neq \pm 1$）：将 $E$ 与 $O$ 以公分母
        $1 - x^2$ 通分并化简。
  \item 展望（每项各用一句话）：指数\omterm{def:g11:func:function}{函数}的偶部分与奇部分
        将在十二年级成为孪生\omterm{def:g11:func:function}{函数}，即双曲余弦与双曲正弦；而
        同样的思想——把对象分解为若干对称的部分，
        再各自按其对称性处理——也支撑着大学分册中
        对声音与信号的分析。
        一般而言，这种分解为何
        如此\emph{有用}？
\end{enumerate}

\medskip
\noindent\textbf{第三部分——变换游戏}
\begin{enumerate}[resume]
  \item 从 $x \mapsto x^2$ 的\omterm{def:g10:functions:graph}{图像}出发，精确描述
        $(x - 3)^2$、\ $x^2 + 2$、\
        $(x - 3)^2 + 2$、\ $-x^2$、\ $2x^2$
        的\omterm{def:g10:functions:graph}{图像}（见\cref{prop:g11:func:transforms}）。
  \item 将 $f(x) = \abs{x - 2} + \abs{x + 2}$ 化为
        在由 $-2$ 与 $2$ 所分出的三个\omterm{def:g10:numbers:interval}{区间}上的分段\omterm{def:g11:func:function}{函数}，
        并描述其\omterm{def:g10:functions:graph}{图像}（它有一段著名的
        “平底”）。
  \item 哪些变换保持偶性？对下列变换判断，并给出证明或反例：
        水平平移 $f(x - 3)$；竖直平移
        $f(x) + 2$；竖直伸缩 $2 f(x)$
        （均作用于\omterm{def:g11:func:parity}{偶函数} $f$）。
  \item 解\omterm{def:g10:algebra:equation}{方程} $\abs{x - 2} + \abs{x + 2} = 6$，
        利用第 12 题中的分段\omterm{def:g11:func:function}{函数}求解。
  \item 逆向思考：一条\omterm{def:g11:quad:parabola}{抛物线}的顶点为 $(2, -3)$
        且经过点 $(0, 1)$。写出该\omterm{def:g11:quad:parabola}{抛物线}的\omterm{def:g10:algebra:equation}{方程}，
        并化为\omterm{thm:g11:quad:canonical}{顶点式}。
\end{enumerate}

\medskip
\noindent\textbf{第四部分——不用微积分求\omterm{def:g10:functions:extrema}{最小值}}
\begin{enumerate}[resume]
  \item 仿照\cref{exo:g11:func:11}，
        对 $g(x) = x + \dfrac4x$，在 $\intoo{0}{+\infty}$ 上
        证明 $g(v) - g(u) = (v - u)\,\dfrac{uv - 4}{uv}$，并确定
        \omterm{def:g10:functions:extrema}{最小值}点，并给出\omterm{def:g10:functions:extrema}{最小值}。
  \item 应用：在所有面积为 $16$ 的矩形中，哪一个的
        周长最小？（记 $P = 2\left(x +
        \frac{16}{x}\right)$，并运用第 16 题中的“机器”。）一位
        老朋友作出结论。
  \item 确定
        $h(x) = \sqrt{x^2 + 1}$ 在整个 $\R$ 上的单调性，
        结合在 $\intco{0}{+\infty}$ 上的复合论证，
        并对其余部分利用奇偶性进行论证。
  \item 最近点问题：求曲线
        $y = \sqrt x$ 上距离 $(3, 0)$ 最近的点。（先使
        \emph{距离的平方}最小——它是关于 $x$ 的\omterm{def:g11:quad:quadratic}{二次函数}——
        再说明为何对平方求最小是合理的。）
  \item 终篇，参观工厂：基本曲线是原材料，
        变换是机器，奇偶性与分解是质量控制，
        差分\omterm{def:g10:algebra:expand}{因式分解}是测量台。每项各用一句话，
        并点出下一章将交付的
        那件利器。
\end{enumerate}
\end{problem}
